% Generated by scripts/assemble.py; corrections belong in source/corrections.json.
\examhead{20}{中国人民大学 2025 年数学分析试卷}{24}
% p025-b19; source PDF page 25
\blockanchor{p025-b19}
\question{1}求最小的实数 $a$, 使得对任意的 $x > 0$, 都有 $\left(1 + \frac{1}{x}\right)^{x + a} > e$.\par
% p025-b20; source PDF page 25
\par\Needspace{6\baselineskip}
\blockanchor{p025-b20}
\question{2}已知 $x_0, y_0, z_0$ 为三个给定的实数, 对正整数 $n \geq 1$, 令\par
\nopagebreak[4]
% p025-b21; source PDF page 25
\blockanchor{p025-b21}
\displayformula{x _ {n} = \frac {y _ {n - 1} + z _ {n - 1}}{2}, y _ {n} = \frac {x _ {n - 1} + z _ {n - 1}}{2}, z _ {n} = \frac {x _ {n - 1} + y _ {n - 1}}{2}.}
% p025-b22; source PDF page 25
\blockanchor{p025-b22}
证明：序列 $\{x_{n}\} ,\{y_{n}\} ,\{z_{n}\}$ 均收敛，并计算它们的极限.\par
% p026-b00; source PDF page 26
\blockanchor{p026-b00}
\question{3}设函数 $f(x)$ 在 $[0,1]$ 上连续, 在 $(0,1)$ 上可导, 且 $f(0) = 0$, $f(1) = 2$. 证明: 存在 $\xi, \eta \in (0,1)$, $\xi \neq \eta$, 满足 $[1 + f'(\xi)][1 + f'(\eta)] = 9$.\par
% p026-b01; source PDF page 26
\blockanchor{p026-b01}
\question{4}设 $n$ 为正整数, 计算极限 $\lim_{n\to \infty}\int_0^1\sqrt[n]{x^n + (1 - x)^n}\mathrm{d}x.$\par
% p026-b02; source PDF page 26
\blockanchor{p026-b02}
\question{5}求三元函数 $f(x_{1}, x_{2}, x_{3}) = \sum_{i=1}^{3} \sum_{j=1}^{3} (e^{x_{i}} - x_{j})^{2}$ 的最小值, 其中 $x_{1}, x_{2}, x_{3}$ 均为实数.\par
% p026-b03; source PDF page 26
\blockanchor{p026-b03}
\question{6}已知定义在 [0,3] 上的非负连续函数 $f(x)$ 满足 $\int_{0}^{3}\frac{1}{1 + f(x)}\mathrm{d}x = 1$，证明：$\int_{0}^{3}\frac{f(x)}{1 + f^2(x)}\mathrm{d}x \leq 1$.\par
% p026-b04; source PDF page 26
\par\Needspace{6\baselineskip}
\blockanchor{p026-b04}
\question{7}设 $D$ 是两条直线 $y = x, y = 4x$ 和两条双曲线 $xy = 1, xy = 4$ 在第一象限围成的区域, 边界 $\partial D$ 的方向为逆时针方向, $F(u)$ 为 $\mathbb{R}$ 上的可导函数, $F'(u) = f(u)$, 且 $f(u)$ 在 $\mathbb{R}$ 中连续, 证明:\par
\nopagebreak[4]
% p026-b05; source PDF page 26
\blockanchor{p026-b05}
\displayformula{\int_ {\partial D} \frac {F (x y)}{y} \mathrm{d} y = \ln 2 \int_ {1} ^ {4} f (u) \mathrm{d} u.}
% p026-b06; source PDF page 26
\blockanchor{p026-b06}
\question{8}设 $\Sigma$ 为上半椭球面 $\frac{x^2}{2} + \frac{y^2}{2} + z^2 = 1, z \geq 0$. 对 $\Sigma$ 上任意一点 $(x, y, z)$, 记 $\rho(x, y, z)$ 为原点到 $\Sigma$ 在 $(x, y, z)$ 处的切平面的距离, 计算 $\iint_{\Sigma} \frac{z}{\rho(x, y, z)} \mathrm{d}S$.\par
% p026-b07; source PDF page 26
\blockanchor{p026-b07}
\question{9}已知对任意的正整数 $n, f_n(x)$ 均为 $[a, b]$ 上的单调函数，且 $\{f_n(x)\}$ 在 $[a, b]$ 上逐点收敛于连续函数 $f(x)$. 证明：$\{f_n(x)\}$ 在 $[a, b]$ 上一致收敛于 $f(x)$.\par
% p026-b08; source PDF page 26
\blockanchor{p026-b08}
\question{10}对正整数 $n$, 判断反常积分 $\int_0^{+\infty} x \sin(x^n) \, \mathrm{d}x$ 是否收敛以及是否绝对收敛, 并说明理由.\par
